\section{Week 2 Packed-Key Prefix}
\label{sec:week2-pack}

\paragraph{Research question and claim.}
\claimid{KG-PREFIX-CALL} asks whether the actual generated mode-2 packer and
KeyGen caller establish
\[
  \prefix_{992}(sk)=vk.
\]
The corresponding paper algorithm is KeyGen in Figure~8 of the pinned
specification \cite{haetae-pdf}. The implementation source is
\pathname{keypair.jazz}, SHA-256
\texttt{8b2ddac2862122d1c9d58767cbc5caa2caad39c8f41fa8096d66ffa77b783e01}.

\paragraph{Procedures and generated source.}
The proof opens the generated procedures
\procname{\_pack\_sk\_m23}, \procname{\_pack\_vec\_eta\_to}, and
\procname{\_pack\_vec2\_eta\_to}. It then lifts their result through
\procname{\_keypair\_full\_m23} and
\procname{crypto\_sign\_keypair\_internal\_mode2\_jazz}.
The imported parent extraction has SHA-256
\texttt{248f8157e348e3f294665d12573a58ee58e860884356bfe82324fda64de5d0b4}.
The independently regenerated packer-focused target has SHA-256
\texttt{7062715ca8fe449f15632c315cafdb316bddad6c39d01fe1f2da51104db6931c}.

\paragraph{Compiled statements.}
Theory \theoryname{ExtractedPackedKeyPrefix} contains:
\begin{verbatim}
hoare [Parent._pack_sk_m23 :
  vkp = vk0 /\ vkbytes = W64.of_int 992 /\
  mcount = W64.of_int 3 /\ kcount = W64.of_int 2
  ==> vk_prefix_eq res vk0 992]
\end{verbatim}
under theorem name
\theoryname{pack\_sk\_m23\_mode2\_vk\_prefix}. It also proves:
\begin{verbatim}
hoare [Parent._keypair_full_m23 :
  k = 2 /\ m = 3 /\ vkbytes = 992
  ==> vk_prefix_eq res.`2 res.`1 992]
\end{verbatim}
as \theoryname{keypair\_full\_m23\_mode2\_return\_prefix}, followed by an
unconditional mode-2 internal-entry lift
\theoryname{keypair\_internal\_mode2\_return\_prefix}.

\paragraph{Proof strategy and invariant.}
The initial copy loop maintains
\[
  \mathsf{copied\_prefix}(skp,vk_0,\mathsf{to\_uint}(i)).
\]
The later-write frame is split rather than assumed:
\begin{itemize}
  \item \theoryname{pack\_vec\_eta\_to\_prefix\_frame} proves that the
  \(3\cdot64\) byte eta-vector region beginning at offset 992 preserves the
  prefix;
  \item \theoryname{pack\_vec2\_eta\_to\_prefix\_frame} does the same for the
  \(2\cdot96\) byte eta2 region; and
  \item \theoryname{mode2\_key\_index\_after\_prefix} proves that the final
  32-byte key write also begins after the prefix.
\end{itemize}

\paragraph{Reachability and non-vacuity.}
This is a Hoare partial-correctness result. It has no hidden termination
premise: every execution that returns satisfies the prefix relation.
Losslessness of the retry loop, the probability of termination, and the output
distribution remain separate obligations. The returned relation is consumed
constructively in
\theoryname{Mode2MuChallengeComposition.keygen\_prefix\_reaches\_mu\_memory};
it is not merely witnessed by arbitrary arrays.

\paragraph{Spec/implementation difference.}
The paper presents structured keys, whereas the Jasmin code lays the packed VK
at the beginning of a 2752-byte SK and appends eta vectors and a 32-byte key.
The theorem validates this concrete layout relation only; it does not yet prove
canonical parsing or equality to the paper's full key distribution.

\paragraph{Verification and status.}
\claimid{KG-PREFIX-CALL} is \status{PROVED} as partial correctness. The target
fresh-compiles with \texttt{-no-eco}; packer regeneration and its output hash
are checked by \pathname{scripts/verify-all.sh}. The next minimum KeyGen
obligations are \claimid{OBL-KG-TERMINATION} and the public-memory lift, not
another array-prefix lemma.
